% ROZDZIAŁ 15 — pełny, z dowodami; nie zakłada twierdzenia Radona–Nikodyma
\chapter[Dualność i geometria przestrzeni $L^p$]{Dualność, wypukłość i słaba zbieżność w przestrzeniach $L^p$}
\label{ch:dualnosc-lp}

\section{Przestrzeń dualna i podstawowy problem reprezentacji}
Niech $(X,\mathcal F,\mu)$ będzie przestrzenią z miarą. Przestrzeń
$L^p(\mu)$ zbudowaliśmy w rozdziale~\ref{ch:przestrzenie-lp} jako
przestrzeń Banacha. Celem jest znalezienie \emph{wszystkich} ciągłych
funkcjonałów liniowych na tej przestrzeni. Należy odróżnić
wykazanie, że pewien wzór daje funkcjonał, od twierdzenia, że
każdy funkcjonał ma tę postać.

\begin{definicja}[Dualność i norma funkcjonału]
\label{def:dual-lp-15}
Dla zespolonej przestrzeni Banacha $E$ jej \emph{przestrzenią dualną}
$E^*$ nazywamy przestrzeń wszystkich ciągłych \emph{zespolono-liniowych}
odwzorowań $\Lambda:E\to\mathbb C$, z normą
$\|\Lambda\|:=\sup\{|\Lambda(f)|:\|f\|_E\le1\}$.
W całym rozdziale parowanie $L^p$ z $L^q$ zapisujemy
$\langle f,g\rangle=\int_X f(x)g(x)\,\mathrm d\mu(x)$;
\emph{nie} wstawiamy tu sprzężenia zespolonego. Jest to parowanie
dwuliniowe, a nie iloczyn skalarny na $L^2$.
\end{definicja}

\begin{lemat}[Izometryczne zanurzenie do dualnej]
\label{lem:canonical-lpq-15}
Niech $1\le p<\infty$ i $q\in(1,\infty]$ będzie wykładnikiem
sprzężonym ($1/p+1/q=1$). Jeśli $p=1$, dodatkowo
zakładamy $\sigma$-skończoność $\mu$. Dla każdego $g\in L^q(\mu)$ wzór
$\Lambda_g(f)=\int_X f(x)g(x)\,\mathrm d\mu(x)$ definiuje element
$(L^p(\mu))^*$ i zachodzi $\|\Lambda_g\|=\|g\|_q$.
Dla $1<p<\infty$ nie zakładamy $\sigma$-skończoności.
\end{lemat}
\begin{proof}
Z nierówności Höldera (twierdzenie~\ref{thm:holder-14}) mamy
$|\Lambda_g(f)|\le\|f\|_p\|g\|_q$, a więc operator jest ciągły.
Dla $1<p<\infty$, jeśli $g\ne0$, definiujemy
$f=\overline{g}|g|^{q-2}$ w punktach, gdzie $g\ne0$, i $f=0$
poza nimi. Gdy $q<2$, zapis w punktach $g=0$ zastępujemy wartością zero;
moduł tak określonej funkcji wynosi $|g|^{q-1}$.
Ponieważ $(q-1)p=q$, mamy $f\in L^p$,
$\Lambda_g(f)=\|g\|_q^q$ i $\|f\|_p=\|g\|_q^{q-1}$.
Wynika stąd nierówność odwrotna $\|\Lambda_g\|\ge\|g\|_q$.

Jeśli $p=1$, kładziemy $M=\|g\|_\infty$. Dla dowolnego
$0<a<M$ zbiór $A=\{|g|>a\}$ ma dodatnią miarę.
Z założonej $\sigma$-skończoności wynika, że istnieje
mierzalny zbiór $B\subseteq A$ o $0<\mu(B)<\infty$:
przecinamy $A$ z jednym ze zbiorów wyczerpania
$X=\bigcup_nX_n$, $\mu(X_n)<\infty$.
Wybierając $f=\chi_B\overline g/(|g|\mu(B))$ na $B$
i $f=0$ poza $B$, otrzymujemy $\|f\|_1=1$ i
$|\Lambda_g(f)|\ge a$. Przechodzimy z $a\uparrow M$.
\end{proof}

\begin{uwaga}[Dlaczego wymagane jest doprecyzowanie dla $p=1$?]
W przestrzeni z miarą, która nie jest półskończona, funkcja $g$
może mieć niezerowy kres istotny, ale całkować się do zera
z każdą funkcją $f\in L^1$. Na przykład na jednoelementowej
przestrzeni $\{*\}$ z miarą $\mu(\{*\})=\infty$ mamy
$L^1=\{0\}$ i $L^\infty\cong\mathbb C$.
Zatem ogólne twierdzenie o izometrycznym zanurzeniu z poprzedniego
lematu należy dla $p=1$ rozumieć tylko przy założeniu
$\sigma$-skończoności. Jest to istotna granica twierdzenia;
przy $1<p<\infty$ problem nie występuje.
\end{uwaga}

\section{Nierówności wypukłości w $L^p$}
W szczególności skonstruujemy reprezentację funkcjonałów
\emph{bez użycia twierdzenia Radona--Nikodyma}. Narzędziem
będzie ścisła wypukłość potęgi normy.

\begin{lemat}[Dwie nierówności dla punktu środkowego]
\label{lem:midpoint-lp-15}
Niech $f,h\in L^p(\mu)$ i $1<p<\infty$. Piszemy
$m=(f+h)/2$ oraz $v=(f-h)/2$.
\begin{enumerate}[label=\textup{(\roman*)}]
\item Jeśli $p\ge2$, to
$\|m\|_p^p+\|v\|_p^p\le(\|f\|_p^p+\|h\|_p^p)/2$.
\item Jeśli $1<p<2$ i $\|f\|_p+\|h\|_p>0$, to
\begin{equation}\label{eq:midpoint-low-15}
\frac{\|f\|_p^p+\|h\|_p^p}{2}-\|m\|_p^p
\ge \frac{p(p-1)}2\,
\frac{\|v\|_p^2}{(\|f\|_p+\|h\|_p)^{2-p}}.
\end{equation}
\end{enumerate}
\end{lemat}
\begin{proof}
Dla $p\ge2$ wystarczy wykazać dla $a,b\in\mathbb C$ nierówność
$|a+b|^p+|a-b|^p\le2^{p-1}(|a|^p+|b|^p)$.
Ponieważ $r=p/2\ge1$ i $u^r+v^r\le(u+v)^r$ dla $u,v\ge0$,
\[
|a+b|^p+|a-b|^p
\le (|a+b|^2+|a-b|^2)^{p/2}
=2^{p/2}(|a|^2+|b|^2)^{p/2}.
\]
Wypukłość $t\mapsto t^{p/2}$ daje
$(|a|^2+|b|^2)^{p/2}\le2^{p/2-1}(|a|^p+|b|^p)$.
Po podzieleniu przez $2^p$ i scałkowaniu otrzymujemy (i).

Dla $1<p<2$ identyfikujemy $\mathbb C$ z $\mathbb R^2$.
Dla $z\ne0$ hesjan funkcji $\psi(z)=|z|^p$ ma postać
$D^2\psi(z)=p|z|^{p-2}I+p(p-2)|z|^{p-4}z\otimes z$.
Najmniejsza wartość własna wynosi
$p(p-1)|z|^{p-2}$. Niech $a,b\in\mathbb C$,
$u=(a+b)/2$, $v=(a-b)/2$ i $S=|a|+|b|$.
Wzdłuż odcinka $t\mapsto u+tv$, $-1\le t\le1$,
zachodzi $|u+tv|\le S$. Dla $v\ne0$ funkcja
$t\mapsto |u+tv|^p$ ma bezwzględnie ciągłą pochodną
na $[-1,1]$. Istotnie, poza ewentualnym pojedynczym
punktem $t_0$, w którym $u+t_0v=0$, jest ona klasy $C^2$,
a jeśli taki punkt istnieje, to druga pochodna jest
w jego otoczeniu oszacowana przez wielokrotność
$|t-t_0|^{p-2}$. Jest to funkcja całkowalna, gdyż
$p-2>-1$; ponadto pierwsza pochodna jest ciągła
również w $t_0$. Można więc zastosować dwukrotnie
wzór Newtona--Leibniza na odcinku, rozbijając go
w razie potrzeby w punkcie $t_0$. Z nierówności
$D^2\psi(z)[v,v]\ge p(p-1)|z|^{p-2}|v|^2$
oraz $p-2<0$ otrzymujemy
\[
\frac{|a|^p+|b|^p}{2}-|u|^p
\ge \frac{p(p-1)}2\,|v|^2S^{p-2}.
\]
W przypadku $S=0$ obie strony rozumiemy jako zero.
Przy $S>0$ wyrażenie z prawej jest nieujemne i mierzalne.
Aby oszacować jego całkę od dołu, używamy Höldera
z wykładnikami $2/p$ i $2/(2-p)$:
\[
\int_X|v|^p\,\mathrm d\mu(x)
\le\left(\int_X |v|^2S^{p-2}\,\mathrm d\mu(x)\right)^{p/2}
\left(\int_XS^p\,\mathrm d\mu(x)\right)^{(2-p)/2}.
\]
Jeżeli pierwsza całka jest nieskończona, teza jest oczywista.
W przeciwnym razie podnosimy obie strony do potęgi $2/p$.
Minkowski zapewnia $\|S\|_p\le\|f\|_p+\|h\|_p$.
Otrzymujemy (ii).
\end{proof}

\begin{twierdzenie}[Jednostajna wypukłość]
\label{thm:uniform-convexity-15}
Dla dowolnej przestrzeni z miarą $(X,\mathcal F,\mu)$
i każdego $1<p<\infty$ przestrzeń $L^p(\mu)$ jest
\emph{jednostajnie wypukła}: dla każdego $\varepsilon>0$
istnieje $\delta>0$ taka, że jeżeli $\|f\|_p=\|h\|_p=1$
i $\|f-h\|_p\ge\varepsilon$, to
$\|(f+h)/2\|_p\le1-\delta$.
\end{twierdzenie}
\begin{proof}
Dla $p\ge2$ lemat~\ref{lem:midpoint-lp-15} daje
$\|(f+h)/2\|_p^p\le1-(\varepsilon/2)^p$.
Dla $1<p<2$ wzór \eqref{eq:midpoint-low-15} daje
$\|(f+h)/2\|_p^p\le1-c_p\varepsilon^2$,
gdzie $c_p=p(p-1)/32>0$ jest dopuszczalną stałą.
Ponieważ $\|f-h\|_p\le2$, dla $0<\varepsilon\le2$
prawa strona jest dodatnia. W obu przypadkach
możemy przyjąć $\delta=1-(1-c\varepsilon^r)^{1/p}>0$
z odpowiednimi $c,r$, jak wyżej; dla $\varepsilon>2$
warunek $\|f-h\|_p\ge\varepsilon$ jest niemożliwy.
\end{proof}

\section{Metoda wariacyjna i dualność dla $1<p<\infty$}
\begin{lemat}[Minimum energii]
\label{lem:variational-lp-15}
Niech $1<p<\infty$ i $\Lambda\in(L^p(\mu))^*$.
Funkcjonał rzeczywisty
$J(f)=\|f\|_p^p/p-\operatorname{Re}\Lambda(f)$
osiąga minimum w pewnym $u\in L^p(\mu)$.
\end{lemat}
\begin{proof}
Z oszacowania $J(f)\ge\|f\|_p^p/p-\|\Lambda\|\|f\|_p$
wynika, że $J$ jest ograniczony z dołu, a każdy ciąg
minimalizujący $(f_n)$ jest ograniczony w $L^p$:
$\|f_n\|_p\le R$ dla pewnego $R<\infty$.
Ponieważ człon liniowy zachowuje średnie, z
lematu~\ref{lem:midpoint-lp-15} wynika przy $p\ge2$
\[
\frac{J(f_n)+J(f_m)}2-J\left(\frac{f_n+f_m}{2}\right)
\ge\frac1p\left\|\frac{f_n-f_m}{2}\right\|_p^p.
\]
Dla $1<p<2$ analogiczna różnica jest co najmniej
$c(p,R)\|f_n-f_m\|_p^2$, gdzie $c(p,R)>0$;
przy $R=0$ ciąg jest zerowy. Niech $a=\inf J$.
Skoro $J((f_n+f_m)/2)\ge a$, a $J(f_n),J(f_m)\to a$,
prawe strony powyższych nierówności zmierzają do zera.
Zatem $(f_n)$ jest ciągiem Cauchy'ego w $L^p$.
Na mocy twierdzenia Riesza--Fischera
\ref{thm:riesz-fischer-14} zbiega do $u\in L^p$.
Ciągłość normy i $\Lambda$ daje $J(u)=a$.
\end{proof}

\begin{twierdzenie}[Dualność $L^p$ dla $1<p<\infty$]
\label{thm:dual-lp-15}
Niech $(X,\mathcal F,\mu)$ będzie \emph{dowolną} przestrzenią
z miarą, $1<p<\infty$, a $q=p/(p-1)$.
Wówczas każdy $\Lambda\in (L^p(\mu))^*$ ma \emph{jedyną}
reprezentację
\begin{equation}\label{eq:dual-lp-15}
\Lambda(f)=\int_X f(x)g(x)\,\mathrm d\mu(x)
\qquad(f\in L^p(\mu))
\end{equation}
przez pewne $g\in L^q(\mu)$. Zachodzi
$\|\Lambda\|=\|g\|_q$. Nie jest potrzebna
$\sigma$-skończoność ani twierdzenie Radona--Nikodyma.
\end{twierdzenie}
\begin{proof}
Z lematu~\ref{lem:variational-lp-15} istnieje minimum $u$
funkcjonału $J(f)=\|f\|_p^p/p-\operatorname{Re}\Lambda(f)$.
Dla $v\in L^p$ i $t\in\mathbb R$ funkcja $t\mapsto J(u+tv)$
ma minimum przy $t=0$. Jej pochodna w zerze wynosi
\[
\int_X |u(x)|^{p-2}\operatorname{Re}(\overline{u(x)}v(x))
\,\mathrm d\mu(x)-\operatorname{Re}\Lambda(v).
\]
Istnienie tej pochodnej uzasadniamy następująco:
dla $|t|\le1$ iloraz różnicowy $|u+tv|^p/p$ jest
oszacowany przez $(|u|+|v|)^{p-1}|v|$,
a funkcja ta jest całkowalna na mocy Höldera,
ponieważ $u,v\in L^p$. Możemy więc zastosować
twierdzenie o zbieżności zmajoryzowanej. Z minimum
wynika $\operatorname{Re}\Lambda(v)=
\operatorname{Re}\int_Xv\overline u|u|^{p-2}\,\mathrm d\mu(x)$.
Zastępując $v$ przez $iv$, otrzymujemy także
równość części urojonych. Zatem \eqref{eq:dual-lp-15}
zachodzi dla $g=\overline u|u|^{p-2}$ (przyjmujemy $g=0$
tam, gdzie $u=0$). Ponieważ $(p-1)q=p$,
$\int_X|g|^q\,\mathrm d\mu(x)=\int_X|u|^p\,\mathrm d\mu(x)<\infty$.

Lemat~\ref{lem:canonical-lpq-15} daje równość norm.
Jeśli $g_1,g_2$ reprezentują ten sam funkcjonał,
to funkcjonał $\Lambda_{g_1-g_2}$ jest zerowy,
a z izometrii wynika $\|g_1-g_2\|_q=0$.
Jest to jednoznaczność w sensie równości prawie wszędzie.
\end{proof}

\begin{przyklad}[Norma funkcjonału całkowego]
Niech $1<p<\infty$, $X=(0,1)$ z miarą Lebesgue'a
oraz $g(x)=x^{-\alpha}$ dla $0<\alpha<1/q$.
Funkcjonał $\Lambda(f)=\int_{(0,1)}f(x)x^{-\alpha}
\,\mathrm d\lambda_1(x)$ należy do $(L^p)^*$, a
$\|\Lambda\|=(1-\alpha q)^{-1/q}$.
Jeśli $\alpha q\ge1$, funkcja $g$ nie należy do $L^q$;
nie istnieje wtedy \emph{ciągłe na całym $L^p$} rozszerzenie
tego funkcjonału całkowego zachowujące tę samą postać.
\end{przyklad}

\section{Przypadek $p=1$ i szczególna rola $L^\infty$}
\begin{twierdzenie}[Dualność $L^1$ przy $\sigma$-skończoności]
\label{thm:dual-l1-15}
Jeżeli $(X,\mathcal F,\mu)$ jest $\sigma$-skończona, to
każdy $\Lambda\in(L^1(\mu))^*$ ma jednoznaczną reprezentację
$\Lambda(f)=\int_X f(x)g(x)\,\mathrm d\mu(x)$
przez $g\in L^\infty(\mu)$, a $\|\Lambda\|=\|g\|_\infty$.
\end{twierdzenie}
\begin{proof}
Wybieramy mierzalne rozłączne $E_n$, których suma jest $X$
i $0\le\mu(E_n)<\infty$. Pomijając zbiory zerowe, kładziemy
$w(x)=\sum_{n\ge1}2^{-n}(1+\mu(E_n))^{-1}\chi_{E_n}(x)$.
Mamy $w>0$ prawie wszędzie i $\int_Xw\,\mathrm d\mu(x)\le1$.
Określamy skończoną miarę $P(A)=\int_Aw\,\mathrm d\mu(x)$.
Dla $h\in L^2(P)$ funkcja $wh$ należy do $L^1(\mu)$, gdyż
$\int_X|wh|\,\mathrm d\mu(x)=\int_X|h|\,\mathrm dP(x)
\le P(X)^{1/2}\|h\|_{L^2(P)}$. Zatem
$T(h):=\Lambda(wh)$ jest ciągłym funkcjonałem na
przestrzeni Hilberta $L^2(P)$. Z twierdzenia o reprezentacji
Riesza dla przestrzeni Hilberta istnieje $g\in L^2(P)$
takie, że $T(h)=\int_X h(x)g(x)\,\mathrm dP(x)$.
Uzasadnienie wykorzystanej reprezentacji jest elementarne:
dla $T\ne0$ wybieramy $u$ prostopadły do $\ker T$ o
$T(u)=1$, po rzutowaniu ortogonalnym na domkniętą
podprzestrzeń $\ker T$; wtedy
$T(h)=\langle h,u\rangle_2/\|u\|_2^2$ przy odpowiedniej
konwencji liniowości. Dla samodzielności przypomnijmy dowód potrzebnego
przypadku twierdzenia o rzucie. Jeśli $K\subset H$ jest
domkniętą podprzestrzenią przestrzeni Hilberta,
a $h\in H$, wybieramy $z_n\in K$ tak, że
$\|h-z_n\|^2\to d^2:=\inf_{z\in K}\|h-z\|^2$.
Tożsamość równoległoboku daje
\[
\|z_n-z_m\|^2
=2\|h-z_n\|^2+2\|h-z_m\|^2
-4\left\|h-\frac{z_n+z_m}{2}\right\|^2
\le2\|h-z_n\|^2+2\|h-z_m\|^2-4d^2.
\]
Stąd $(z_n)$ jest ciągiem Cauchy'ego i ma granicę
$z\in K$. Dla każdego $k\in K$ funkcja
$t\mapsto\|h-z-tk\|^2$ ma minimum przy $t=0$;
biorąc również $tk$ z czysto urojonym współczynnikiem,
otrzymujemy $h-z\perp K$. Tak otrzymujemy potrzebny rzut
bez odwołania do zewnętrznego twierdzenia.

Dla każdego $h\in L^2(P)\cap L^1(P)$ mamy ponadto
$|T(h)|\le\|\Lambda\|\int_X|h|\,\mathrm dP(x)$.
W szczególności, podstawiając
$h=\chi_A\overline g/|g|$ (równe zero przy $g=0$),
otrzymujemy
$\int_A|g|\,\mathrm dP(x)\le\|\Lambda\|P(A)$
dla dowolnego $A\in\mathcal F$. Stąd
$|g|\le\|\Lambda\|$ $P$-prawie wszędzie; ponieważ
$w>0$ $\mu$-prawie wszędzie, także $\mu$-prawie wszędzie.

Jeśli $f\in L^1(\mu)$, to $h=f/w\in L^1(P)$
(zera $w$ pomijamy). Funkcje
$h_n=h\chi_{\{|h|\le n\}}$ należą do
$L^2(P)$, a $wh_n\to f$ w $L^1(\mu)$, co wynika
z dominacji przez $|f|$. Zatem
$\Lambda(f)=\lim_nT(h_n)=\lim_n\int_Xh_ng\,\mathrm dP(x)
=\int_Xfg\,\mathrm d\mu(x)$.
Izometrię i jednoznaczność otrzymujemy z
lematu~\ref{lem:canonical-lpq-15} w wariancie $p=1$
i z $\sigma$-skończonością.
\end{proof}

\begin{uwaga}[Dlaczego dualna $L^\infty$ jest większa?]
Nie jest prawdą, że $(L^\infty)^*=L^1$ nawet dla
miary liczącej na $\mathbb N$. Wtedy $L^\infty=\ell^\infty$.
Na podprzestrzeni $c$ ciągów zbieżnych funkcjonał
$T((x_n))=\lim_nx_n$ jest ciągły i ma normę $1$.
Twierdzenie Hahna--Banacha przedłuża go do
$\widetilde T\in(\ell^\infty)^*$ o tej samej normie.
Jeżeli byłby postaci $\widetilde T(x)=\sum_n a_nx_n$
z $(a_n)\in\ell^1$, to dla wektorów jednostkowych $e_n$
mielibyśmy $a_n=\widetilde T(e_n)=0$ (każdy $e_n\in c_0$),
lecz $\widetilde T(\chi_X)=1$. Sprzeczność pokazuje,
że taki funkcjonał nie pochodzi z $\ell^1$.
W rozdziałach o miarach ze znakiem i Radonie--Nikodymie
poznamy dodatkowo reprezentacje przez miary, lecz nie
należy utożsamiać dualnej $L^\infty$ z miarami
\emph{przeliczalnie} addytywnymi.
\end{uwaga}

\begin{wniosek}[Dokładne maksimum w nierówności Höldera]
\label{cor:holder-sharp-15}
Niech $1<p<\infty$, $q=p/(p-1)$, a
$g\in L^q(\mu)$, $g\neq0$. Istnieje
$f\in L^p(\mu)$ o $\|f\|_p=1$ takie, że
\[
\int_X f(x)g(x)\,\mathrm d\mu(x)=\|g\|_q.
\]
Można wybrać
$f(x)=\overline{g(x)}|g(x)|^{q-2}/\|g\|_q^{q-1}$
przy $g(x)\neq0$ oraz $f(x)=0$ przy $g(x)=0$.
\end{wniosek}
\begin{proof}
Mamy $|f|^p=|g|^q/\|g\|_q^q$,
ponieważ $(q-1)p=q$, zatem $\|f\|_p=1$.
Ponadto $fg=|g|^q/\|g\|_q^{q-1}$, więc
jego całka wynosi $\|g\|_q$. Te rachunki
pokazują, że stała $1$ w nierówności Höldera
jest optymalna niezależnie od przestrzeni z miarą.
\end{proof}

\section{Refleksywność, zbieżność słaba i słaba-*}
\begin{definicja}[Refleksywność i zbieżność słaba]
\label{def:weak-lp-15}
Dla przestrzeni Banacha $E$ zanurzenie kanoniczne
$J:E\to E^{**}$ jest dane przez $(Jx)(\Lambda)=\Lambda(x)$.
Przestrzeń $E$ jest \emph{refleksywna}, jeżeli $J$ jest
surjekcją. Ciąg $x_n$ \emph{zbiega słabo} do $x$,
co zapisujemy $x_n\rightharpoonup x$, jeżeli
$\Lambda(x_n)\to\Lambda(x)$ dla każdego $\Lambda\in E^*$.
Na dualnej $E^*$ mówimy o \emph{zbieżności słabej-*}
$\Lambda_n\overset{*}{\rightharpoonup}\Lambda$,
gdy $\Lambda_n(x)\to\Lambda(x)$ dla każdego $x\in E$.
Są to dwie na ogół różne topologie na $E^*$.
\end{definicja}

\begin{twierdzenie}[Refleksywność $L^p$]
\label{thm:reflexive-lp-15}
Dla dowolnej przestrzeni z miarą $(X,\mathcal F,\mu)$
i $1<p<\infty$ przestrzeń $L^p(\mu)$ jest refleksywna.
\end{twierdzenie}
\begin{proof}
Niech $q=p/(p-1)$. Z twierdzenia~\ref{thm:dual-lp-15}
izometrycznie identyfikujemy $(L^p)^*$ z $L^q$ przez
$g\mapsto\Lambda_g$, gdzie $\Lambda_g(f)=\int_Xfg\,\mathrm d\mu(x)$.
Ponieważ $1<q<\infty$, to samo twierdzenie zastosowane
z zamienionymi rolami $p,q$ utożsamia $(L^q)^*$ z $L^p$.
Pod tymi identyfikacjami kanoniczny obraz $J(f)$
działa na $g\in L^q$ jako $J(f)(g)=\int_Xfg\,\mathrm d\mu(x)$.
Każdy element $(L^q)^*$ ma tę postać, co dowodzi
surjektywności $J$.
\end{proof}

\begin{twierdzenie}[Kryterium słabej zbieżności w $L^p$]
\label{thm:weak-lp-15}
Niech $1<p<\infty$, $q=p/(p-1)$ i $f_n,f\in L^p(\mu)$.
Wtedy $f_n\rightharpoonup f$ w $L^p$ dokładnie wtedy, gdy
\[
\int_Xf_n(x)g(x)\,\mathrm d\mu(x)
\longrightarrow\int_Xf(x)g(x)\,\mathrm d\mu(x)
\quad\text{dla każdego }g\in L^q(\mu).
\]
Jeśli $f_n\to f$ w normie $L^p$, to $f_n\rightharpoonup f$.
\end{twierdzenie}
\begin{proof}
Pierwsza równoważność jest dokładnie połączeniem
definicji~\ref{def:weak-lp-15} z
reprezentacją~\ref{thm:dual-lp-15}.
Druga wynika z $|\Lambda(f_n-f)|\le\|\Lambda\|\|f_n-f\|_p$.
\end{proof}

\begin{twierdzenie}[Dolna półciągłość normy]
\label{thm:weak-lsc-15}
W dowolnej przestrzeni Banacha $E$, jeśli $x_n\rightharpoonup x$,
to $\|x\|\le\liminf_n\|x_n\|$.
\end{twierdzenie}
\begin{proof}
Dla każdego $\Lambda\in E^*$ z $\|\Lambda\|\le1$
mamy $|\Lambda(x)|=\lim_n|\Lambda(x_n)|
\le\liminf_n\|x_n\|$.
Na mocy twierdzenia Hahna--Banacha istnieje
funkcjonał normujący wektor $x$, co daje
$\|x\|\le\liminf_n\|x_n\|$.
W przypadku $E=L^p$, $1<p<\infty$, można uniknąć
Hahna--Banacha: z lematu~\ref{lem:canonical-lpq-15}
$\|x\|_p=\sup_{\|g\|_q\le1}|\int_Xxg\,\mathrm d\mu(x)|$.
\end{proof}

\begin{przyklad}[Słaba zbieżność bez zbieżności normowej]
Na $X=(0,1)$ z miarą Lebesgue'a niech
$r_n(x)=\operatorname{sign}(\sin(2^{n+1}\pi x))$,
przy dowolnym ustaleniu wartości na zbiorze miejsc zerowych.
Funkcja $r_n$ przyjmuje naprzemiennie wartości $1$ i $-1$
na kolejnych przedziałach diadycznych długości $2^{-n-1}$.
Dla każdego $1<p<\infty$ mamy $\|r_n\|_p=1$, lecz
$r_n\rightharpoonup0$ w $L^p(0,1)$. Istotnie, całka
$r_n$ po dowolnym przedziale diadycznym ustalonego
rzędu jest zerowa dla wszystkich dostatecznie dużych $n$.
W konsekwencji $\int r_ng\to0$ dla funkcji schodkowych
na przedziałach diadycznych. Ich gęstość w $L^q(0,1)$
i Hölder pozwalają przenieść zbieżność na każde $g\in L^q$.
\end{przyklad}

\begin{twierdzenie}[Podciąg słabo zbieżny: wariant separowalny]
\label{thm:weak-subseq-15}
Niech $1<p<\infty$, a $L^q(\mu)$ będzie separowalna.
Wtedy z każdego ciągu ograniczonego $(f_n)$ w $L^p(\mu)$
można wybrać podciąg słabo zbieżny do pewnego $f\in L^p(\mu)$.
\end{twierdzenie}
\begin{proof}
Niech $(g_j)_{j\ge1}$ będzie przeliczalnym gęstym
zbiorem w $L^q$. Dla każdego $j$ ciąg liczb
$\int_Xf_n g_j\,\mathrm d\mu(x)$ jest ograniczony przez
$M\|g_j\|_q$, gdzie $M=\sup_n\|f_n\|_p$.
Kolejno wybieramy podciągi, na których te liczby
są zbieżne, i przyjmujemy podciąg diagonalny
$(f_{n_k})$. Granica definiuje funkcjonał liniowy na
liniowej powłoce $(g_j)$, ograniczony przez $M\|g\|_q$.
Przedłużamy go przez ciągłość na całe $L^q$.
Z twierdzenia~\ref{thm:dual-lp-15} zastosowanego do $q$
funkcjonał ten ma postać $g\mapsto\int_Xfg\,\mathrm d\mu(x)$
dla pewnego $f\in L^p$. Gęstość i jednolite
oszacowanie $M$ zapewniają zbieżność całek
$\int_Xf_{n_k}g\to\int_Xfg$ dla wszystkich $g\in L^q$.
Kryterium~\ref{thm:weak-lp-15} daje tezę.
\end{proof}

\begin{uwaga}[Zastosowanie i ograniczenie]
W ogólnym przypadku refleksywności przestrzeni Banacha
odpowiednikiem ostatniego twierdzenia jest twierdzenie
Eberleina--\v Smuliana (zwartość słaba i zwartość ciągowa).
Nie dowodzimy go tutaj: nie jest twierdzeniem czysto
miarowym. Dowód diagonalny powyżej jest samodzielny
i wystarcza między innymi dla $L^p(\mathbb R^d)$,
$1<p<\infty$.
\end{uwaga}

\section{Zastosowanie: minimalizacja funkcjonałów całkowych}
\begin{twierdzenie}[Minimalizacja wypukłej energii]
\label{thm:minimize-lp-15}
Niech $1<p<\infty$, $g\in L^q(\mu)$ i $\alpha>0$.
Funkcjonał $\mathcal E(f)=\alpha\|f\|_p^p/p
-\operatorname{Re}\int_X f g\,\mathrm d\mu(x)$
ma dokładnie jedno minimum w $L^p(\mu)$. Jego
punkt minimalny jest dany prawie wszędzie przez
\[
f_*(x)=\alpha^{-1/(p-1)}\,\overline{g(x)}\,|g(x)|^{q-2},
\]
przy umowie $f_*(x)=0$ dla $g(x)=0$.
\end{twierdzenie}
\begin{proof}
Istnienie minimum dla $\alpha>0$ wynika z dowodu
lematu~\ref{lem:variational-lp-15} (stały dodatni
współczynnik nie zmienia argumentu). Funkcja
$z\mapsto|z|^p$ jest ściśle wypukła na $\mathbb C$,
a więc całka jest ściśle wypukła na klasach $L^p$;
stąd jednoznaczność. Warunek pierwszego rzędu daje
$\alpha|f_*|^{p-2}\overline{f_*}=g$.
Rozwiązanie tego równania w punktach $g\ne0$
ma wskazaną postać, ponieważ $q-1=1/(p-1)$;
$|f_*|^p=\alpha^{-q}|g|^q$ jest całkowalne.
To dowodzi zarówno istnienia wskazanej funkcji,
jak i spełnienia warunku minimalności (wypukłość
$\mathcal E$ sprawia, że warunek jest wystarczający).
\end{proof}

\section{Mapa zależności}
Dowiedliśmy, że dla $1<p<\infty$ kolejno zachodzą:
nierówności punktu środkowego, jednostajna wypukłość,
istnienie minimizatora, reprezentacja dualna i
refleksywność. Przypadek $p=1$ wykorzystał zamiast tego
reprezentację Hilbertowską po zmianie miary na skończoną.
Nie zakładaliśmy rozkładu Radona--Nikodyma; w rozdziale 18
zobaczymy jego szerszą wersję dla samych miar.
